\chapter{Systematic Macro}\label{ch:s2:systematic-macro}

Value, momentum and carry, measured country by country in bonds, currencies and equity indices, make a macro portfolio that no economist designed.
Asness, Moskowitz and Pedersen found value and momentum premia in eight markets and asset classes. Each correlated more with its counterparts in other
asset classes than the asset classes did with each other, and value and momentum were negatively correlated. On this chapter's synthetic market of
thirty futures, value earns a Sharpe ratio of 0.84, momentum 0.78, carry 0.33 and an \omterm{def:s2:systematic-macro:surprise}{economic-surprise signal} 0.29. Nearly uncorrelated, they combine
to 1.15. Measure value from prices instead of from a fundamental anchor and it becomes momentum's mirror image, $-0.84$, and the combination falls to
0.31. The build is \texttt{firm.sysmacro}.

\section{Signals across countries and assets}

The universe is Book~8's synthetic futures (chapter~19): ten equity indices, ten government bonds and ten currencies, one of each per country, over
thirty years, with a persistent drift (the trend) and a carry premium. \texttt{firm.sysmacro} adds two planted effects. Each market carries a
valuation gap from fair value with a standard deviation of 10\% and a half-life of three years; the price includes the gap, and its reversion is the
value premium. Each country has an economic-surprise index with a half-life of 60 days that moves its markets over the following days.

\begin{definition}[Macro value signal]\label{def:s2:systematic-macro:value}
A \emph{macro value signal}\index{macro value signal} ranks countries' assets by their price against a fundamental anchor, such as purchasing-power
parity for a currency, a real-yield or term-premium model for a bond, or earnings and book value for an equity index, buying the cheap and selling the
dear.
\end{definition}

Each signal is traded long-short within each asset class, reset monthly by rank, and scaled to 10\% volatility on its trailing quarter
(\cref{lst:s2:systematic-macro:signals}). Value uses an anchor that each market's analyst gets wrong by a persistent 10\%. Momentum is the
twelve-month return skipping the last month (Book~8, chapter~19), and carry is Book~8's carry (chapter~20). The price-based variant of value is minus
the last five years' spot move, the return less the carry earned.

\section{Macro-data signals}

\begin{definition}[Macro momentum]\label{def:s2:systematic-macro:momentum}
\emph{Macro momentum}\index{macro momentum} takes long positions in the assets of countries whose fundamental economic trends (growth, inflation,
external balance, monetary policy) are improving and short positions where they are deteriorating, the macro-data counterpart of price momentum.
\end{definition}

\begin{definition}[Economic-surprise signal]\label{def:s2:systematic-macro:surprise}
An \emph{economic-surprise signal}\index{economic-surprise signal} measures how much a country's recent data releases have exceeded or fallen short
of forecasts, and positions its assets for the markets' delayed response: equities and the currency up and bonds down after positive surprises.
\end{definition}

Brooks describes \omterm{def:s2:systematic-macro:momentum}{macro momentum} as a systematic global macro strategy that is long assets whose macroeconomic trends are improving and short those
whose trends are deteriorating, with low correlation to traditional asset classes and diversification in equity bear markets and rising-yield
environments. The synthetic surprise signal is a plain version of it. Each country's index moves its equity index up by 4\% a year per unit, its
currency up by 3\%, and its bond down by 3\%, and the trader sees the index with noise of half its size.

\section{Combining signals}

\begin{center}
\small
\begin{tabular}{@{}lccc@{}}
\toprule
24.5 years, 10\% vol each & Sharpe & max drawdown & corr.\ with momentum\\
\midrule
value (anchor) & 0.84 & $-25\%$ & $-0.03$\\
value from prices (reversal) & $-0.84$ & & $-0.43$\\
momentum & 0.78 & $-23\%$ & 1\\
carry & 0.33 & $-32\%$ & 0.04\\
surprise & 0.29 & $-38\%$ & $-0.01$\\
\midrule
all four & 1.15 & $-21\%$ &\\
all four, value from prices & 0.31 & &\\
all but value & 0.81 & &\\
\bottomrule
\end{tabular}
\end{center}

Four nearly uncorrelated signals combine as the arithmetic says. Equal risk budgets give a Sharpe ratio of about the sum of the four divided by the
square root of four, 1.12, against 1.15 measured (\cref{fig:s2:systematic-macro:cumulative}). The combination's worst drawdown, 21\%, is smaller than
any single signal's. Value earns its place: without it the combination falls to 0.81.

How value is measured decides whether it helps. In this market the trend lasts about a year, and the last five years' price move still carries some of
it. Minus that move is short the markets that trended up, so the price-based value signal becomes the mirror of momentum: $-0.84$ against 0.78, a
correlation of $-0.43$. Combined, the two cancel, and the four-signal book earns only 0.31. Asness and co-authors found value and momentum both
positive and negatively correlated. A synthetic market reproduces that only if value's anchor is not built from the same prices momentum uses.

\begin{omfigure}
\begin{tikzpicture}
\begin{axis}[width=11cm,height=5.6cm,xlabel={\small year},ylabel={\small cumulative return (\%, simple sum)},tick label style={font=\footnotesize},
  xmin=0,xmax=24.5,ymin=-20,ymax=300,grid=major,grid style={gray!25},table/col sep=comma,
  legend style={font=\scriptsize,at={(0.5,-0.3)},anchor=north},legend columns=5]
\addplot[omDef,thick] table[x=year,y=value]{figdata/strategies-2/16-systematic-macro/cumulative.csv};
\addplot[omThm,thick] table[x=year,y=momentum]{figdata/strategies-2/16-systematic-macro/cumulative.csv};
\addplot[gray,thick] table[x=year,y=carry]{figdata/strategies-2/16-systematic-macro/cumulative.csv};
\addplot[gray!60,thick,dashed] table[x=year,y=surprise]{figdata/strategies-2/16-systematic-macro/cumulative.csv};
\addplot[black,very thick] table[x=year,y=combined]{figdata/strategies-2/16-systematic-macro/cumulative.csv};
\legend{value,momentum,carry,surprise,combined}
\end{axis}
\end{tikzpicture}

\omcaption{Four synthetic macro signals, each scaled to 10\% volatility and traded long-short within equities, bonds and currencies, and their
equal-risk combination (also at 10\%), cumulative sums of daily returns over 24.5 years. Data: \texttt{s2\_sysmacro.books}.}\label{fig:s2:systematic-macro:cumulative}
\end{omfigure}

\section{What the record shows}

The public record supports the pattern, not the numbers. Asness and co-authors found value and momentum premia everywhere they looked, with a common
factor structure across asset classes and a partial link to global funding liquidity. Brooks's \omterm{def:s2:systematic-macro:momentum}{macro momentum} adds the trends in economic data. The
chapter's Sharpe ratios are those of a planted market (\cref{fig:s2:systematic-macro:sharpe}). They are there to show how signals combine, not to
forecast what a real book would earn. What carries over is the method: measure each signal on its own, measure their correlations, give each a risk
budget, and check that value and momentum are not measured from the same data.

\begin{omfigure}
\begin{tikzpicture}
\begin{axis}[width=11.5cm,height=5.6cm,ybar,bar width=12pt,ylabel={\small Sharpe ratio},tick label style={font=\footnotesize},
  xmin=0.4,xmax=8.6,ymin=-1.3,ymax=1.4,xtick={1,2,3,4,5,6,7,8},
  xticklabels={value,reversal,momentum,carry,surprise,{all four},{all four, price value},{all but value}},
  x tick label style={font=\scriptsize,rotate=25,anchor=north east},
  grid=major,grid style={gray!25},table/col sep=comma,nodes near coords,every node near coord/.append style={font=\scriptsize}]
\addplot[fill=omThm!70,draw=omThm] table[x=x,y=sr]{figdata/strategies-2/16-systematic-macro/sharpe.csv};
\end{axis}
\end{tikzpicture}

\omcaption{Sharpe ratios of the synthetic macro signals and three combinations: all four (value from the anchor), all four with value from prices,
and three without value. Data: \texttt{s2\_sysmacro.stats}.}\label{fig:s2:systematic-macro:sharpe}
\end{omfigure}

\section{Strategy files}

\begin{strategyfile}{Cross-country value}\label{strat:s2:systematic-macro:value}
\sfield{Who pays you, and why} Investors who extrapolate, and a fundamental anchor that prices return to over years.
\sfield{Instruments and venues} Equity index, bond and currency futures and forwards.
\sfield{Signal} Price against a fundamental anchor: purchasing-power parity, real yields, earnings yields.
\sfield{Sizing and execution} Long-short within asset class; monthly; risk-budgeted.
\sfield{Costs} Futures rolls; low turnover.
\sfield{How it dies} Anchors that are wrong for years; trends that run further than value can wait.
\sfield{Horizon, capacity, infrastructure} Years; large capacity.
\sfield{Backtest honestly} Anchors as they could be computed then, not revised data.
\sfield{Sources} Asness, Moskowitz and Pedersen (2013); this chapter: 0.84 from an anchor, $-0.84$ from prices.
\end{strategyfile}

\begin{strategyfile}{Macro momentum}\label{strat:s2:systematic-macro:momentum}
\sfield{Who pays you, and why} Markets that price economic news slowly.
\sfield{Instruments and venues} As above.
\sfield{Signal} Trends in growth, inflation, external balances and policy; price momentum as a companion.
\sfield{Sizing and execution} Long-short; monthly.
\sfield{Costs} Moderate turnover.
\sfield{How it dies} Turning points; data revisions.
\sfield{Horizon, capacity, infrastructure} Months; point-in-time macro data.
\sfield{Backtest honestly} First-release data, not revised.
\sfield{Sources} Brooks (2017); this chapter: price momentum 0.78.
\end{strategyfile}

\begin{strategyfile}{Carry across asset classes}\label{strat:s2:systematic-macro:carry}
\sfield{Who pays you, and why} As in Book~8, chapter~20: holders of low-carry assets, and the crash risk of high carry.
\sfield{Instruments and venues} Futures and forwards in three asset classes.
\sfield{Signal} Each market's carry.
\sfield{Sizing and execution} Long-short within class; risk-budgeted.
\sfield{Costs} Rolls.
\sfield{How it dies} Carry crashes.
\sfield{Horizon, capacity, infrastructure} Months.
\sfield{Backtest honestly} Carry from futures curves, not assumed rates.
\sfield{Sources} Book~8, chapter~20; this chapter: 0.33.
\end{strategyfile}

\begin{strategyfile}{Economic-surprise tilt}\label{strat:s2:systematic-macro:surprise}
\sfield{Who pays you, and why} Investors who react slowly to data surprises.
\sfield{Instruments and venues} Country equity, bond and currency futures.
\sfield{Signal} Recent releases against consensus forecasts, by country.
\sfield{Sizing and execution} Long-short; rebalanced as data arrive.
\sfield{Costs} Higher turnover than value.
\sfield{How it dies} Forecasters and markets catch up faster.
\sfield{Horizon, capacity, infrastructure} Weeks to months; a data calendar with forecasts.
\sfield{Backtest honestly} Release times and first prints.
\sfield{Sources} This chapter: 0.29; no public performance figure verified.
\end{strategyfile}

\begin{strategyfile}{Combined macro factor portfolio}\label{strat:s2:systematic-macro:combined}
\sfield{Who pays you, and why} The sum of the above.
\sfield{Instruments and venues} All of the above.
\sfield{Signal} Value, momentum, carry and macro data, with equal or estimated risk budgets.
\sfield{Sizing and execution} Target volatility at the portfolio level.
\sfield{Costs} Netting reduces turnover.
\sfield{How it dies} Correlations that rise together in a crisis.
\sfield{Horizon, capacity, infrastructure} Months; a multi-asset futures book.
\sfield{Backtest honestly} Signals defined before the test; correlations measured out of sample.
\sfield{Sources} Asness, Moskowitz and Pedersen (2013); this chapter: 1.15 combined.
\end{strategyfile}

\section{Tutorial: designed by no one}\label{tut:s2:systematic-macro:tut}

\begin{tutorial}
\textbf{Goal.} Plant value and macro-surprise effects in the synthetic futures, build five signals, and combine them. \textbf{End state:} the table
and the two figures.
\begin{tutsteps}
\item \textbf{Signals and books}.
\omcode{code/firm/sysmacro/firm_sysmacro.py}{75}{114}{The five signals, the within-class long-short weights and the vol-targeted book.}\label{lst:s2:systematic-macro:signals}
\item \textbf{The combinations}.
\omcode{code/strategies-2/16-systematic-macro/python/s2_sysmacro.py}{30}{51}{Each signal, three combinations, and their statistics.}\label{lst:s2:systematic-macro:books}
\item \textbf{Run} \texttt{stats()}, \texttt{correlations()} and \texttt{fig\_sysmacro.py}.
\end{tutsteps}
\textbf{What to change next.} Estimate risk budgets from trailing correlations instead of equal budgets; add a slow-moving anchor error that drifts;
test the price-based value on a market whose trend lasts three months.
\end{tutorial}

\section{Build: systematic macro}\label{bld:s2:systematic-macro:sysmacro}

\begin{build}
\bfield{Purpose} Cross-asset value, momentum, carry and macro-surprise signals, and their risk-budgeted combination.
\bfield{Interface} \texttt{MacroConfig(\dots)}, \texttt{macro\_market(cfg)}, \texttt{signal(m, name, t)}, \texttt{portfolio(m, name, cfg)},
\texttt{combine(books, cfg)}.
\bfield{Rules} Signals from data up to the day before; long-short within each class; each book scaled on its own trailing volatility.
\bfield{Acceptance tests} \texttt{code/firm/sysmacro/tests/}: within-class weights are neutral with unit gross; the market's shapes and signals; the
combination targets its volatility.
\bfield{Stretch} Estimated risk budgets; real macro data; transaction costs.
\end{build}

\begin{omsources}
\item C.~S.~Asness, T.~J.~Moskowitz and L.~H.~Pedersen, ``Value and momentum everywhere'', \emph{Journal of Finance} 68(3), 2013.
\item J.~Brooks, ``A half century of macro momentum'', AQR white paper, 2017.
\end{omsources}

\section{Exercises}

\begin{exercise}[$\star$]\label{exo:s2:systematic-macro:1}
Four uncorrelated signals have Sharpe ratios of 0.84, 0.78, 0.33 and 0.29 and equal risk budgets. What Sharpe ratio does their combination have?
\end{exercise}

\begin{exercise}[$\star$]\label{exo:s2:systematic-macro:2}
Two signals have Sharpe ratios of 0.84 and 0.78 and a correlation of $-0.43$. What is the Sharpe ratio of an equal-risk combination?
\end{exercise}

\begin{exercise}[$\star$]\label{exo:s2:systematic-macro:3}
A valuation gap has a half-life of three years. What share of it closes in one year?
\end{exercise}

\begin{exercise}[$\star\star$]\label{exo:s2:systematic-macro:4}
Why does value measured from five-year price moves become momentum's mirror image in this market?
\end{exercise}

\begin{exercise}[$\star\star$]\label{exo:s2:systematic-macro:5}
Why must a macro-data backtest use first-release data?
\end{exercise}

\begin{exercise}[$\star\star$]\label{exo:s2:systematic-macro:6}
Why trade signals long-short within each asset class rather than across them?
\end{exercise}

\begin{exercise}[$\star\star\star$]\label{exo:s2:systematic-macro:7}
\emph{Coding.} Rerun \texttt{macro\_market} with \texttt{MacroConfig(anchor\_noise=0.3)}. What happens to value, and to the combination?
\end{exercise}

\begin{exercise}[$\star\star\star$]\label{exo:s2:systematic-macro:8}
\emph{Find the flaw.} ``Our combination has a Sharpe ratio of 1.15, so we will run it at 30\% volatility.''
\end{exercise}

\section{Problem: Designed by No One}

\begin{problem}[{Weekend problem --- a systematic macro book}]\label{pb:s2:systematic-macro:1}
The chapter's synthetic futures and the public record.

\medskip\noindent\textbf{Part I --- Signals.}
\begin{enumerate}
\item Define the \omterm{def:s2:systematic-macro:value}{macro value signal} and give three anchors.
\item Define \omterm{def:s2:systematic-macro:momentum}{macro momentum} and the \omterm{def:s2:systematic-macro:surprise}{economic-surprise signal}.
\item How is each signal traded and scaled?
\item What did Asness, Moskowitz and Pedersen find?
\end{enumerate}

\medskip\noindent\textbf{Part II --- The planted market.}
\begin{enumerate}[resume]
\item Describe the valuation gap and the surprise index.
\item How does the analyst see each?
\item Why does the anchor's error matter?
\item What did Brooks describe?
\end{enumerate}

\medskip\noindent\textbf{Part III --- Results.}
\begin{enumerate}[resume]
\item Give each signal's Sharpe ratio.
\item Give the correlations.
\item Give the three combinations' Sharpe ratios.
\item Why does value from prices fail here?
\end{enumerate}

\medskip\noindent\textbf{Part IV --- The verdict.}
\begin{enumerate}[resume]
\item State the \emph{named result}: each signal's Sharpe ratio and the combination's, with the correlations that make it work.
\item What does the arithmetic of uncorrelated signals predict?
\item What happens to the combination in a crisis?
\item How would you set risk budgets?
\item How would you backtest the book honestly?
\item Which strategy file has the highest turnover?
\item How does this chapter relate to Book~8's trend and carry chapters?
\item In one sentence: why does a combination of simple signals work?
\end{enumerate}
\end{problem}

\section{Interview questions}

\begin{interviewq}[$\star$ \iqroles{researcher}]\label{iq:s2:systematic-macro:1}
How would you measure value in currencies, bonds and equity indices?
\end{interviewq}

\begin{interviewq}[$\star\star$ \iqroles{researcher}]\label{iq:s2:systematic-macro:2}
Why are value and momentum negatively correlated, and why is that useful?
\end{interviewq}

\begin{interviewq}[$\star\star$ \iqroles{trader}]\label{iq:s2:systematic-macro:3}
Your macro book's momentum signal is short Japanese bonds and its value signal is long. What do you do?
\end{interviewq}

\begin{interviewq}[$\star\star$ \iqroles{risk}]\label{iq:s2:systematic-macro:4}
How would you stress a systematic macro book?
\end{interviewq}

\begin{interviewq}[$\star\star$ \iqroles{developer}]\label{iq:s2:systematic-macro:5}
Design a point-in-time macro database for backtests.
\end{interviewq}

\begin{interviewq}[$\star\star\star$ \iqroles{researcher}]\label{iq:s2:systematic-macro:6}
Show that equal-risk combinations of $n$ signals with Sharpe ratios $S_i$ and common pairwise correlation $\rho$ have a Sharpe ratio of $\sum_i S_i /
\sqrt{n + n(n - 1)\rho}$.
\end{interviewq}
